% GENERATED FILE. Edit canonical lessons, then run npm run generate:books.
% canonical-source-hash: fnv1a64:291a5b0a34468dd4
% canonical-lessons: UR-C55-alu, UR-C55-tamatar, UR-C55-gajar, UR-C55-chini, UR-C55-namak

\chapter{Potato, Tomato, Carrot, Sugar, Salt}
\label{ch:ur-potato-tomato-carrot-sugar-salt}
\hlchaptermodality{55}

\begin{chapteropening}
\textbf{By the end of this chapter:} \emph{I can say a potato, a tomato, a carrot, sugar and salt.}

Complete the last lesson of chapter 55: I can say a potato, a tomato, a carrot, sugar and salt.
\end{chapteropening}

\section[\texorpdfstring{\ur{آلو} (ālū)}{آلو (ālū)}]{\ur{آلو} (ālū) — a potato}

\label{lesson:UR-C55-alu}



\begin{quote}
Before the new one: say the Urdu for an apple, then the Urdu for a banana.
\end{quote}

\subsection*{You'll want to know: \ur{آلو}}
\textbf{\ur{آلو}} — \emph{ālū} — \textquotedblleft{}a potato\textquotedblright{}.

\begin{cousinweb}
\textbf{\ur{آلو}} (\emph{ālū}) is a tadbhava, a word inherited from Sanskrit and reshaped by sound change. It comes from \emph{ālu}, \textquotedblleft{}an edible root\textquotedblright{} (CDIAL 1388). When the potato came to India, it took the name of an older root vegetable.
\end{cousinweb}

\begin{grammarlens}[title={What you've built}]
One more word for this chapter.
\end{grammarlens}

\subsection*{Your turn}
\begin{itemize}
\raggedright
  \item \emph{Say it:} \emph{ālū}
  \item \emph{Say it:} \emph{ālū}, once more
  \item \emph{Say it:} say \emph{ālū}
  \item \emph{Recall:} say the Urdu for an apple, then the Urdu for a banana, then say \emph{ālū} again
\end{itemize}

\begin{tcolorbox}[breakable,colback=teal!4,colframe=teal!35!black,title={Before you move on}]
What does \ur{آلو} mean? (\textquotedblleft{}A potato\textquotedblright{}.) Say it once more.
\end{tcolorbox}
\section[\texorpdfstring{\ur{ٹماٹر} (ṭamāṭar)}{ٹماٹر (ṭamāṭar)}]{\ur{ٹماٹر} (ṭamāṭar) — a tomato}

\label{lesson:UR-C55-tamatar}



\begin{quote}
Before the new one: say the Urdu for a banana, then the Urdu for a potato.
\end{quote}

\subsection*{You'll want to know: \ur{ٹماٹر}}
\textbf{\ur{ٹماٹر}} — \emph{ṭamāṭar} — \textquotedblleft{}a tomato\textquotedblright{}.

\begin{cousinweb}
\textbf{\ur{ٹماٹر}} (\emph{ṭamāṭar}) is a loan from English \emph{tomato}.
\end{cousinweb}

\begin{grammarlens}[title={What you've built}]
One more word for this chapter.
\end{grammarlens}

\subsection*{Your turn}
\begin{itemize}
\raggedright
  \item \emph{Say it:} \emph{ṭamāṭar}
  \item \emph{Say it:} \emph{ṭamāṭar}, once more
  \item \emph{Say it:} say \emph{ṭamāṭar}
  \item \emph{Recall:} say the Urdu for a banana, then the Urdu for a potato, then say \emph{ṭamāṭar} again
\end{itemize}

\begin{tcolorbox}[breakable,colback=teal!4,colframe=teal!35!black,title={Before you move on}]
What does \ur{ٹماٹر} mean? (\textquotedblleft{}A tomato\textquotedblright{}.) Say it once more.
\end{tcolorbox}
\section[\texorpdfstring{\ur{گاجر} (gājar)}{گاجر (gājar)}]{\ur{گاجر} (gājar) — a carrot}

\label{lesson:UR-C55-gajar}



\begin{quote}
Before the new one: say the Urdu for a potato, then the Urdu for a tomato.
\end{quote}

\subsection*{You'll want to know: \ur{گاجر}}
\textbf{\ur{گاجر}} — \emph{gājar} — \textquotedblleft{}a carrot\textquotedblright{}.

\begin{cousinweb}
\textbf{\ur{گاجر}} (\emph{gājar}) is a tadbhava, a word inherited from Sanskrit and reshaped by sound change. It comes from \emph{gārjara}, \textquotedblleft{}carrot\textquotedblright{} (CDIAL 4140). The Sanskrit word is known only from the old word-lists.
\end{cousinweb}

\begin{grammarlens}[title={What you've built}]
One more word for this chapter.
\end{grammarlens}

\subsection*{Your turn}
\begin{itemize}
\raggedright
  \item \emph{Say it:} \emph{gājar}
  \item \emph{Say it:} \emph{gājar}, once more
  \item \emph{Say it:} say \emph{gājar}
  \item \emph{Recall:} say the Urdu for a potato, then the Urdu for a tomato, then say \emph{gājar} again
\end{itemize}

\begin{tcolorbox}[breakable,colback=teal!4,colframe=teal!35!black,title={Before you move on}]
What does \ur{گاجر} mean? (\textquotedblleft{}A carrot\textquotedblright{}.) Say it once more.
\end{tcolorbox}
\section[\texorpdfstring{\ur{چینی} (chīnī)}{چینی (chīnī)}]{\ur{چینی} (chīnī) — sugar}

\label{lesson:UR-C55-chini}



\begin{quote}
Before the new one: say the Urdu for a tomato, then the Urdu for a carrot.
\end{quote}

\subsection*{You'll want to know: \ur{چینی}}
\textbf{\ur{چینی}} — \emph{chīnī} — \textquotedblleft{}sugar\textquotedblright{}.

\begin{cousinweb}
\textbf{\ur{چینی}} (\emph{chīnī}) means \textquotedblleft{}Chinese\textquotedblright{}: sugar was named after China. CDIAL lists Hindi \emph{cīnā}, \textquotedblleft{}Chinese\textquotedblright{}, under Sanskrit \emph{cīna}, \textquotedblleft{}Chinese\textquotedblright{} (CDIAL 4841); \emph{chīnī} is the same name with the ending \emph{-ī}.
\end{cousinweb}

\begin{grammarlens}[title={What you've built}]
One more word for this chapter.
\end{grammarlens}

\subsection*{Your turn}
\begin{itemize}
\raggedright
  \item \emph{Say it:} \emph{chīnī}
  \item \emph{Say it:} \emph{chīnī}, once more
  \item \emph{Say it:} say \emph{chīnī}
  \item \emph{Recall:} say the Urdu for a tomato, then the Urdu for a carrot, then say \emph{chīnī} again
\end{itemize}

\begin{tcolorbox}[breakable,colback=teal!4,colframe=teal!35!black,title={Before you move on}]
What does \ur{چینی} mean? (\textquotedblleft{}Sugar\textquotedblright{}.) Say it once more.
\end{tcolorbox}
\section[\texorpdfstring{\ur{نمک} (namak)}{نمک (namak)}]{\ur{نمک} (namak) — salt}

\label{lesson:UR-C55-namak}



\begin{quote}
Before the new one: say the Urdu for a carrot, then the Urdu for sugar.
\end{quote}

\subsection*{You'll want to know: \ur{نمک}}
\textbf{\ur{نمک}} — \emph{namak} — \textquotedblleft{}salt\textquotedblright{}.

\begin{cousinweb}
\textbf{\ur{نمک}} (\emph{namak}) is a loan, usually traced to Persian \emph{namak}, \textquotedblleft{}salt\textquotedblright{}.
\end{cousinweb}

\begin{grammarlens}[title={What you've built}]
One more word for this chapter.
\end{grammarlens}

\subsection*{Your turn}
\begin{itemize}
\raggedright
  \item \emph{Say it:} \emph{namak}
  \item \emph{Say it:} \emph{namak}, once more
  \item \emph{Say it:} say \emph{namak}
  \item \emph{Recall:} say the Urdu for a carrot, then the Urdu for sugar, then say \emph{namak} again
\end{itemize}

\begin{tcolorbox}[breakable,colback=teal!4,colframe=teal!35!black,title={Before you move on}]
What does \ur{نمک} mean? (\textquotedblleft{}Salt\textquotedblright{}.) Say it once more.
\end{tcolorbox}
